\chapter{Justificativa}
\label{cap:justificativa}

A pesquisa se justifica pela crescente digitalização das experiências humanas e a
interconexão cada vez maior entre o físico e o virtual, que demandam uma compreensão
aprofundada de como as pessoas percebem e interagem com os mais diversos ambientes e
contextos. Este trabalho se insere nesse cenário, buscando explorar como a tecnologia pode não
apenas mediar, mas também auxiliar na compreensão e no enriquecimento da relação entre
indivíduos, lugares e suas identidades. A seguir, detalham-se os pilares que sustentam a relevância e
a oportunidade desta pesquisa.

\section{Relevância dos Princípios ESG no Contexto Atual}
\label{sec:relevancia-esg}

Embora o foco principal desta pesquisa não seja a avaliação direta do desempenho \gls{ESG} de
organizações -- sigla que representa os pilares Ambiental (\textit{Environmental}), Social (\textit{Social}) e de
Governança (\textit{Governance}) e se refere às práticas voltadas à sustentabilidade ambiental,
responsabilidade social e à qualidade de sua administração e transparência -- a consideração desses
princípios é fundamental no desenvolvimento e na aplicação de tecnologias que interagem com a
sociedade e o ambiente. Ferramentas que promovem o entendimento da identidade local e a
valorização do patrimônio cultural, por exemplo, dialogam com as dimensões sociais e de
governança ao fomentar o bem-estar comunitário, o respeito à diversidade e o uso consciente dos
recursos.

\section{A Crescente Necessidade de Compreender a Percepção e Identidade de Lugar}
\label{sec:necessidade-percepcao}

A forma como indivíduos e comunidades percebem e se identificam com os lugares que
habitam ou visitam é fundamental para o seu bem-estar, coesão social e para a preservação do
patrimônio cultural e natural. Na era digital, ferramentas interativas, como aplicativos de navegação
e plataformas de experiência turística e urbana, desempenham um papel cada vez mais significativo
na mediação dessa relação. Compreender profundamente como essas tecnologias moldam a
percepção dos usuários sobre os espaços, e como contribuem para a formação ou o reforço de sua
identidade territorial, é crucial. Este entendimento permite o desenvolvimento de intervenções
digitais mais sensíveis e eficazes, capazes de promover um sentimento de pertencimento genuíno e
uma valorização mais profunda dos contextos locais, indo além da simples funcionalidade para
tocar em aspectos emocionais e identitários da experiência humana.

\section{O Potencial da Análise de Sentimento e da Design Science para Inovação e Impacto}
\label{sec:potencial-analise}

Para desvendar as percepções e os sentimentos subjetivos dos usuários em relação aos
lugares, a análise de sentimento -- compreendida como a área do Processamento de Linguagem
Natural (\gls{NLP}) que visa extrair e identificar os estados afetivos e as opiniões subjetivas presentes
em textos \cite{liu2023esg} -- emerge como uma ferramenta analítica de grande potencial para
capturar as percepções públicas e acadêmicas. Ao aplicar algoritmos capazes de interpretar opiniões
e emoções expressas em dados textuais gerados a partir da interação com ferramentas digitais, é
possível obter resultados sobre a experiência do usuário que dificilmente seriam capturados por
métodos tradicionais. A condução desta pesquisa e o desenvolvimento da solução tecnológica
associada serão pautados pela metodologia da Design Science \cite{wieringa2014}. Esta
abordagem é particularmente adequada por focar na criação e avaliação de artefatos que resolvem
problemas identificados no mundo real, ao mesmo tempo em que geram conhecimento científico. A
combinação da análise de sentimento com os preceitos da Design Science justifica-se pela
oportunidade de criar uma ferramenta inovadora e validada para melhor identificar a percepção e a
identidade de lugar, contribuindo tanto para o avanço tecnológico quanto para a promoção de
experiências mais ricas, humanizadas e potencialmente mais alinhadas com um desenvolvimento
local sustentável.
